\begin{proof}
	Fix a vertex \( v \) and an edge \( e \) leaving \( v \). For each Eulerian tour \( \cE \) that
	starts with \( v \) and \( e \) and
	each vertex \( u\ne v \), consider the edge of first arrival at \( u \).
	\begin{claim}
		The edges of first arrivals with \( e \) forms an arborescence \( T \) at \( v \).
	\end{claim}
	Let us order all incoming edge at every \( u\ne v \) in the order of traversal in \( \cE \), these
	order have following property:
	\begin{itemize}
		\item Among all the incoming edges at \( u\ne v \), the first edge belongs to \( T \).
		\item At \( v \), the edge \( e \) is the last in the ordering.
	\end{itemize}

	\begin{lemma}
		There is a bijection between:
		$$
			\begin{gathered}
				\{\text{Eulerian tours ending with } e \text{ at } v\} \\
				\updownarrow \\
				\left\{
				\begin{aligned}
					 & \text{Arborescence } T \text{ at } v, \text{ and orderings on} \\
					 & \text{incoming edges satisfying the above conditions}
				\end{aligned}
				\right\}
			\end{gathered}
		$$
	\end{lemma}

	\begin{proof}
		Suppose we have
		\begin{itemize}
			\item Spanning arb at \( v \).
			\item An ordering on incoming edges satisfying the above two bullet points.
		\end{itemize}
		Start with \( v \) and traverse \( e \) which the edge that enters \( v \) last. Traverse the
		graph by visiting the larger edge at each stop \( u\ne v \). The only time we get stuck is when
		we visit \( v \) the last time (balanced).
		\begin{claim}
			At this point, we have traversed all edges in \( D \).
		\end{claim}
		\begin{proof}[\textbf{Proof of Claim}]
			Suppose otherwise, we left some edges out. Among the vertices incident to these edges, choose
			\( u \) to be the closest to \( v \), otherwise have \( u\ne v \). Since \( D \) is balanced,
			we have the same number of incoming edges and out comming edges at \( u \), there is at least
			one incoming edge at \( u \) not traversed yet. By condition \( 1 \), we have not traversed
			the unique edge in \( T \) at \( u \). Let \( u' \) be the vertex adjacent to \( e' \), with
			\( u' \xrightarrow{e'} u \) by definition of \( u' \), \( u' \) must be strictly closer to \(
			v\) that \( u \), leading to a contradiction $\lightning$.
		\end{proof}
	\end{proof}
\end{proof}

\begin{definition}[\textbf{Directed Laplacian}]
	Let \( D \) be a digraph with vertex set \( V = \{v_1,\ldots, v_p\} \) and with \( l_i \)
	loops at vertex \( v_i \). Let \( \LL(D) \) be the \underline{directed laplacian} defined by:
	$$
		L_{ij} = \begin{cases}
			-m_{ij},                  & \text{if } i \neq j \text{ and there are } m_{ij} \text{ edges with} \\
			                          & \text{initial vertex } v_i \text{ and final vertex } v_j             \\ \\
			\text{outdeg}(v_i) - l_i, & \text{if } i = j.
		\end{cases}
	$$
\end{definition}

\begin{theorem}
	\( \det(\LL_{0}) = \tau(D,v_p)\)
\end{theorem}

